\documentclass[12pt]{article}
\usepackage{mathtools}
\usepackage{amssymb}
\usepackage{fontspec}
\setmainfont{EB Garamond}

% Set page size and margins
\usepackage[a4paper,top=2cm,bottom=1.5cm,left=1cm,right=1cm]{geometry}
\usepackage{fancyhdr}
\pagestyle{fancy}
\setlength{\parindent}{0pt}

% Useful packages
\usepackage{minted}
\usepackage{graphicx}
\graphicspath{{./media/}}
\usepackage{subfig}
\usepackage{float}
\usepackage[colorlinks=true, allcolors=blue]{hyperref}
\usepackage{siunitx}


\title{\vspace{-2cm} Θεμελειώσεις Κρυπτογραφίας \\
        Homework 1}
\author{Μαρινόπουλος Χριστόφορος 10522}
\date{Εαρινό εξάμηνο 2025}

\lhead{Homework 1}
\chead{Θεμελειώσεις Κρυπτογραφίας}
\rhead{Μαρινόπουλος Χριστόφορος}
\begin{document}
\maketitle

\section*{Εκφώνηση}
Έστω το Eλληνικό αλφάβητο {α,β,γ,…,ω}. Δίνουμε διαδοχικά τιμές
α:0,β:1,…,ω:24. To κρυπτοσύστημα μετατόπισης π.χ. με κλειδί k=3, ορίζεται
E(k,α)=(0+3)mod 24 = 3→δ. Δηλ. Αν το κλειδί είναι k με 0<=k<=23, τότε
Ε(k,X)=(αριθμητική τιμή του X + k)mod 24, ενώ η αποκρυπτογράφηση
D(k,Y)=(αριθμητική τιμή του Y - k)mod 24
Αν έχουμε να κρυπτογραφήσουμε μια λέξη εφαρμόζουμε την κρυπτογράφηση σε κάθε γράμμα
ξεχωριστά.

\subsection*{Ερώτημα 1}
Έστω ότι έχουμε την παρακάτω κατανομή επί του χώρου μηνυμάτων
Pr[M=’α’]=0.7, Pr[M=’ω’]=0.3.
Να υπολογιστεί η πιθανότητα Pr[C=’β’].

\subsection*{Απάντηση}
Ορίζω το σύνολο:
\begin{equation*}
L = \{ x \in \mathbb{N} \mid 1 \leq x \leq 23 \}
\end{equation*}
\newline
Για να κρυπτογραφηθεί το μήνυμα α:0 στο μήνυμα β:1 πρέπει να υπάρχει $ k \, \epsilon \, L $ τέτοιο ώστε: 
$$
(0 + \text{k})mod24 = 1 \Rightarrow k = 1  
$$
Αντίστοιχα, για να κρυπτογραφηθεί το μήνυμα ω:23 στο μήνυμα β:1 πρέπει  να υπάρχει $ k \, \epsilon \, L $ τέτοιο ώστε:
$$
(23 + \text{k})mod 24 = 1 \Rightarrow k = 2 
$$
Άρα, το γράμμα β μπορεί να προκύψει εάν το αρχικό γράμμα ήταν είτε α είτε ω, με κλειδιά k=1 και k=2 αντίστοιχα. 

Συνεπώς η πιθανότητα Pr[C = \text{β}] είναι:
$$
Pr[C = \text{β}] = Pr[M = \text{α}] * Pr[k = 1] + Pr[M = \text{ω}] * Pr[k = 2]
$$
Θεωρώντας πως η κατανομή Κ είναι ομοιόμορφη προκύπτει:
$$
Pr[k = 1] = Pr[k = 2] = \frac{1}{24}
$$
Με αντικατάσταση καταλήγουμε πως:
$$
Pr[C = \text{β}] =0.7 * \frac{1}{24} + 0.3 * \frac{1}{24} = \frac{1}{24}
$$
\hfill \blacksquare
\newpage
\subsection*{Ερώτημα 2}
Έστω ότι έχουμε την παρακάτω κατανομή επί του χώρου μηνυμάτων
Pr[M=’ενα’]=0.3, Pr[M=’δυο’]=0.7.
Να υπολογιστεί η πιθανότητα Pr[C=’θπδ’].

\subsection*{Απάντηση}
Ορίζω το σύνολο:
\begin{equation*}
L = \{ x \in \mathbb{N} \mid 1 \leq x \leq 23 \}
\end{equation*}
\newline
Για να κρυπτογραφηθεί το μήνυμα ενα στο μήνυμα θπδ πρέπει να υπάρχει $ k \, \epsilon \, L $ τέτοιο ώστε να ικανοποιούνται οι ακόλουθες σχέσεις ταυτόχρονα:

\begin{equation*}
\begin{array}{ccc}
\begin{aligned}
(4 + k) \mod 24 &= 7 \\
(12 + k) \mod 24 &= 15 \\
(0 + k) \mod 24 &= 3
\end{aligned}
& \Rightarrow &
k = 3
\end{array}
\end{equation*}
Για να κρυπτογραφηθεί το μήνυμα δύο στο μήνυμα θπδ πρέπει να υπάρχει $ k \, \epsilon \, L $ τέτοιο ώστε να ικανοποιούνται οι ακόλουθες σχέσεις ταυτόχρονα:

\begin{equation*} \label{eq:2}
\begin{aligned}
(3 + k) \mod 24 &= 7 \\
(19 + k) \mod 24 &= 15 \\
(14 + k) \mod 24 &= 3
\end{aligned}
\end{equation*}
Προφανώς δεν υπάρχει $ k \, \epsilon \, L $ τέτοιο ώστε να ικανοποιούνται ταυτόχρονα οι παραπάνω σχέσεις. Επομένως με αυτή την μέθοδο κρυπτογράφησης το μήνυμα δυο δεν μπορεί να κρυπτογραφηθεί στο θπδ.
\newline
\newline
Θεωρώντας πως η κατανομή Κ είναι ομοιόμορφη προκύπτει:
$$
Pr[k = 3] = \frac{1}{24}
$$
Συνεπώς:
$$
Pr[C = \text{θπδ}] = Pr[M = \text{ενα}] * Pr[k = 3] = 0.3 * \frac{1}{24} = 1/80
$$
\hfill \blacksquare
\end{document}
